\subsection{A quantitative consequence for rational Diophantine problems}
\label{subsec:quantitative-mrdp}

There is a useful refinement of the reduction in
Section~\ref{sec:reduction}. The literature gives a bound of $12N$
witnesses for computably enumerable subsets of $\Q$ when $\Z$ has
a definition with $N$ witnesses \cite[Remark 8.23]{DDF2025}.
Using a small positivity condition together with the constructive
intersection bound gives the following alternative.

\begin{lemma}\label{lem:integer-positivity-rational}
For an integer $t$,
\[
 t\ge0\quad\Longleftrightarrow\quad
 \exists u,v,w\in\Q:\ t=u^2+v^2+w^2+w.
\]
\end{lemma}
\begin{proof}
If $t\ge0$, the three-square theorem applies to $4t+1$.
Its three squares have one odd and two even entries, since
$4t+1\equiv1\pmod4$. Write them as $(2w+1)^2,(2u)^2,(2v)^2$
with integers $u,v,w$ and divide by four. Conversely, a rational
solution gives
$t+1/4=u^2+v^2+(w+1/2)^2\ge0$. The assumption $t\in\Z$
then gives $t\ge0$. This is the integer positivity identity used
in \cite[Equation (1.1)]{Sun2021}, with its witnesses now allowed
to be rational. The integer hypothesis must be retained.
\end{proof}

\begin{proposition}\label{prop:quantitative-mrdp}
Suppose $\Z$ is defined in $\Q$ with $N\ge1$ existential witnesses.
Every computably enumerable subset of $\Q$ then has a one-polynomial
definition with at most $11N+1$ existential witnesses. Furthermore,
polynomial solvability over $\Q$ restricted to at most $9N+2$
variables is undecidable.
\end{proposition}
\begin{proof}
We use the published integer representation theorem
\cite[Theorem 1.1(i)]{Sun2021}: for every computably enumerable
$B\subseteq\mathbb Z_{\ge0}$ there is an integer polynomial $g_B$
such that, for $c\ge0$,
\[
 c\in B\quad\Longleftrightarrow\quad
 \exists y_1,\ldots,y_9\in\Z:\
 y_9\ge0,\quad g_B(c,y_1,\ldots,y_9)=0.
\]
For a computably enumerable $A\subseteq\Q$, encode integer pairs
by the injective polynomial
\[
 q(a)=2a^2-a,\qquad
 f(a,b)=(q(a)+q(b))(q(a)+q(b)+1)+2q(b).
\]
The map $q:\Z\to\mathbb Z_{\ge0}$ is injective, and the second
formula is twice the usual pairing polynomial on nonnegative integers.
For a fixed sum of the two arguments its values are distinct, and
successive sums have disjoint ranges, proving injectivity.
The set
\[
 B=\{f(a,b):a,b\in\Z,\ b\ne0,\ a/b\in A\}
\]
is computably enumerable and contained in $\mathbb Z_{\ge0}$.
Thus $x\in A$ is equivalent to the existence of
$y_0,\ldots,y_9\in\Q$ satisfying
\[
 xy_0,y_0,y_1,\ldots,y_9\in\Z,\quad
 y_0\ne0,\quad y_9\ge0,\quad
 g_B(f(xy_0,y_0),y_1,\ldots,y_9)=0.
\]
There are eleven integrality predicates, each using $N$ private
witnesses. Use Lemma~\ref{lem:integer-positivity-rational} for
$y_9\ge0$, writing $y_9=u^2+v^2+w^2+w$, and use
$y_0\eta-1=0$ for $y_0\ne0$.

Regard $g_B(f(xy_0,y_0),y_1,\ldots,y_8,Y)$ as a polynomial in $Y$.
On the closed stratum where all its coefficients vanish, impose
$y_9=0$; on the complement retain its root equation. This preserves
the final projection to $x$. Indeed, on the former stratum the only
other restrictions on $y_9$ are integrality and nonnegativity, both
satisfied by zero. Its private integrality witnesses can be chosen
afresh and we may set $u=v=w=0$. We do not claim preservation of
the projection to every larger tuple of private coordinates.

Normalize each integrality polynomial in its last private witness as
in the proof of Proposition~\ref{prop:intersection-rank}. Retain
$y_0,\ldots,y_8$, the first $N-1$ witnesses of each integrality
predicate, and $u,v$. Besides $x$, this retains
\[
 9+11(N-1)+2=11N
\]
coordinates. Over every geometric retained tuple, $y_9$ has finitely
many choices; then each last integrality witness has finitely many
choices, $w$ satisfies the monic quadratic
$w^2+w+u^2+v^2-y_9=0$, and $\eta$ is unique whenever it exists.
The resulting constructible system has finite geometric fibers.
Theorem~\ref{thm:finite-fibre-one-variable} replaces all these
remaining choices by one witness, giving $11N+1$ in total.

For the last assertion fix a nonrecursive computably enumerable
$B\subseteq\mathbb Z_{\ge0}$ and use $g_B$ directly, with its
input $c$ supplied as a nonnegative integer parameter. Only the nine
integrality predicates on $y_1,\ldots,y_9$ and the three-witness
positivity predicate are needed. Apply the same normalization to
$g_B(c,y_1,\ldots,y_8,Y)$. Retaining $y_1,\ldots,y_8$, the first
$N-1$ private coordinates in each predicate, and $u,v$ leaves
$8+9(N-1)+2=9N+1$ coordinates. One final finite-choice witness
gives $9N+2$. An algorithm deciding rational solvability in this
many variables would
therefore decide $c\in B$, a contradiction. This step concerns an
integer input parameter; it makes no assertion that the same formula
defines $B$ on arbitrary rational inputs without an integrality test
on that parameter.
\end{proof}

With $N=42$, the conditional 43-variable construction therefore gives
$463$ existential witnesses for each computably enumerable subset of
$\Q$, and undecidability of rational polynomial solvability in at most
$380$ variables. These conclusions have the same arithmetic assumptions
as that construction. Using only the source-preserving 558-variable
construction instead, $N=557$ gives the bounds $6128$ and $5015$.
These are quantitative consequences of the defining-polynomial method,
not claims of optimality among all undecidability reductions. In
particular a specialized undecidability theorem and a definition of the
ordinary subset $\Z\subset\Q$ are different objectives. The first
bound improves $12N$ when $N>1$. No polynomial achieving these
MRDP-derived counts has been expanded in this report.
